\documentclass[11pt]{article}

\usepackage[T1]{fontenc}
\usepackage{amsmath,amssymb,amsthm}
\usepackage{hyperref}

\title{Consistency of Normal-Theory FIML under Non-normality and
Missingness}
\author{}
\date{\today}

\newcommand{\E}{\mathbb{E}}
\newcommand{\tr}{\operatorname{tr}}
\newcommand{\avar}{\operatorname{avar}}
\newcommand{\vech}{\operatorname{vech}}
\newcommand{\Cov}{\operatorname{Cov}}
\newcommand{\Var}{\operatorname{Var}}
\newcommand{\argmin}{\operatorname*{arg\,min}}
\newcommand{\argmax}{\operatorname*{arg\,max}}
\newcommand{\plim}{\operatorname*{plim}}

\theoremstyle{plain}
\newtheorem{theorem}{Theorem}
\newtheorem{lemma}{Lemma}
\newtheorem{corollary}{Corollary}
\theoremstyle{remark}
\newtheorem*{remark}{Remark}

\begin{document}
\maketitle

\section*{What is being proved}

Two questions, kept separate because their answers differ.

\begin{enumerate}
  \item \textbf{Non-normality.} Is the normal-theory (Gaussian) FIML estimator
        consistent for the mean and covariance when the data are non-normal?
        \emph{Yes, unconditionally}, for complete data and under MCAR, assuming
        only finite second moments and a correctly specified structure. The
        argument is elementary and uses no distributional assumption at all.
  \item \textbf{MAR.} Is it consistent under data missing at random?
        \emph{Yes, but conditionally.} Consistency holds if and only if, for the
        incomplete variables, the conditional means given the always-observed
        causes are linear and the conditional covariances are homoscedastic.
        Marginal non-normality is harmless; conditional non-linearity or
        heteroscedasticity is not.
\end{enumerate}

The second answer reconciles the textbook statement ``FIML is consistent under
MAR'' (Rubin's ignorability \cite{rubin1976}; Arminger and Sobel
\cite{armingersobel1990}; Yuan and Bentler \cite{yuanbentler2000}) with the
folklore that it is not, and with simulations in which the MAR bias looks like a
small-sample effect that shrinks with $N$. Both camps are right about different
regimes; the precise dividing line is Theorem~\ref{thm:mar} below.

Throughout, $X_1,\ldots,X_N$ are i.i.d.\ $p$-vectors with finite second moments,
true mean $\mu_0$ and covariance $\Sigma_0\succ0$. The saturated parameter is
$\eta=(\mu,\vech\Sigma)$ with true value $\eta_0$. A structural model writes
$\eta=\eta(\theta)$. ``Consistent'' means convergence in probability to the
full-data target as $N\to\infty$.

\section{The saturated Gaussian MLE is the sample moments}

The whole non-normality question collapses once we notice that with complete
data the saturated Gaussian likelihood has a closed-form maximizer that never
saw a normal assumption.

\begin{lemma}\label{lem:sat}
With complete data, the maximizer of the Gaussian log-likelihood
$\sum_i\big[-\tfrac12\log|\Sigma|-\tfrac12(x_i-\mu)'\Sigma^{-1}(x_i-\mu)\big]$
over $(\mu,\Sigma\succ0)$ is
\[
  \hat\mu = \bar x = \tfrac1N\textstyle\sum_i x_i,
  \qquad
  \hat\Sigma = S = \tfrac1N\textstyle\sum_i (x_i-\bar x)(x_i-\bar x)'.
\]
\end{lemma}

\begin{proof}
The mean score $\sum_i\Sigma^{-1}(x_i-\mu)=0$ gives $\hat\mu=\bar x$ for any
$\Sigma\succ0$. Substituting and maximizing over $\Sigma$ is the standard
Gaussian profile problem whose solution is the empirical second-moment matrix
$S$.
\end{proof}

So the saturated FIML point estimate \emph{is} the method-of-moments estimator
of $(\mu,\Sigma)$. It is a Gaussian \emph{quasi}-likelihood when the data are
non-normal (the general misspecified-likelihood theory of White
\cite{white1982}), but its location is fixed by the first two empirical moments
alone.

\begin{theorem}[Complete-data consistency, distribution-free]\label{thm:complete}
If $\E\lVert X\rVert^2<\infty$ then $\hat\eta\xrightarrow{p}\eta_0$. If in
addition the structure is correctly specified, with $\eta(\theta_0)=\eta_0$ for a
unique identified $\theta_0$, and the discrepancy
$\theta\mapsto F\{\hat\eta;\eta(\theta)\}$ is continuous with a unique minimizer,
then $\hat\theta\xrightarrow{p}\theta_0$.
\end{theorem}

\begin{proof}
By the law of large numbers $\bar x\xrightarrow{p}\mu_0$ and
$S\xrightarrow{p}\Sigma_0$, so $\hat\eta\xrightarrow{p}\eta_0$ by
Lemma~\ref{lem:sat}. The fitting function $F$ is minimized at the truth and is
continuous in its first argument, so the standard argmin (Wald) consistency
theorem, together with the continuous mapping theorem, gives
$\hat\theta\xrightarrow{p}\theta_0$. Normality is used nowhere.
\end{proof}

\begin{remark}
This is consistency of the \emph{point estimate}. Validity of normal-theory
standard errors and of the unscaled likelihood-ratio reference law is a strictly
stronger property that does require either asymptotic-robustness conditions
\cite{browne1984,andersonamemiya1988} or a sandwich correction
\cite{white1982}. Theorem~\ref{thm:complete} is exactly why the point estimates
survive non-normality while the test statistics need the FMG/MLR machinery of
the companion note.
\end{remark}

\section{MCAR: still distribution-free}

With missing data the saturated MLE loses the closed form of
Lemma~\ref{lem:sat} (it needs EM), so we work directly with the FIML score and
show it is an \emph{unbiased estimating equation} at the full-data moments. This
makes the $h_i$ argument fully explicit.

\paragraph{The FIML likelihood and its score.}
Let unit $i$ have observed index set $o_i\subseteq\{1,\dots,p\}$ and observed
subvector $x_{i,o_i}$. FIML contributes the Gaussian \emph{marginal} log-density
of that subvector,
\[
  \ell_i(\eta)
   = -\tfrac12\log\det\Sigma_{o_io_i}
     -\tfrac12\,(x_{i,o_i}-\mu_{o_i})'\Sigma_{o_io_i}^{-1}(x_{i,o_i}-\mu_{o_i}),
\]
dropping the constant $-\tfrac12|o_i|\log2\pi$. Write the residual
$r_i=x_{i,o_i}-\mu_{o_i}$ and the observed-block precision
$W_i=\Sigma_{o_io_i}^{-1}$. Differentiating in the saturated coordinates
$\eta=(\mu,\vech\Sigma)$ gives the casewise score
$h_i(\eta)=\partial\ell_i/\partial\eta$, with a mean block and a covariance
block,
\[
  \frac{\partial\ell_i}{\partial\mu_{o_i}} = W_i\,r_i,
  \qquad
  \frac{\partial\ell_i}{\partial\Sigma_{o_io_i}}
    = \tfrac12\big(W_i\,r_i r_i'\,W_i - W_i\big),
\]
and \emph{zero} in every coordinate indexing an unobserved mean $\mu_j$
($j\notin o_i$) or a covariance entry that touches an unobserved variable. (The
covariance block is read in the symmetric matrix and mapped to $\vech$ by the
duplication matrix, which has full column rank and so does not move the root.)
Thus $h_i$ is a \emph{sparse} vector in $\eta$-space, supported exactly on the
moments pattern $o_i$ observes; FIML is the stacked estimating equation
$\sum_i h_i(\eta)=0$, in which each moment's row aggregates only the units that
see it.

\paragraph{Unbiasedness at the truth.}
Read off the two blocks: the mean block has mean zero and the covariance block
has mean zero exactly when
\[
  \E[r_i]=0
  \qquad\text{and}\qquad
  \E[r_i r_i'] = \Sigma_{0,o_io_i},
\]
i.e.\ when the observed subvector carries its \emph{population} first two
moments. (For the covariance block, $\E[W_i r_i r_i' W_i - W_i]
= W_i(\E[r_i r_i'] - \Sigma_{0,o_io_i})W_i = 0$ at $\eta_0$.) That is the only
distributional fact the argument needs, and it is exactly what MCAR supplies.

\begin{theorem}[MCAR consistency, distribution-free]\label{thm:mcar}
Suppose missingness is MCAR, the response pattern $R$ independent of $X$, with
finite second moments, and every entry of $(\mu_0,\Sigma_0)$ identified by some
observed pattern. Then $\hat\eta\xrightarrow{p}\eta_0$; with a correctly
specified structure, $\hat\theta\xrightarrow{p}\theta_0$.
\end{theorem}

\begin{proof}
Condition on the realized pattern. Under MCAR, $R\perp X$, so conditioning on
$R_i=o$ does not change the law of $X_{i,o}$: it keeps the population marginal,
whence $\E[r_i\mid R_i=o]=0$ and $\E[r_i r_i'\mid R_i=o]=\Sigma_{0,oo}$ at
$\eta_0$. By the display above both score blocks then have conditional
expectation zero, and the unobserved coordinates of $h_i$ are identically zero,
so $\E\{h_i(\eta_0)\mid R_i=o\}=0$ for every pattern $o$. Averaging over the
pattern distribution gives $\E\{h_i(\eta_0)\}=0$: the full-data $\eta_0$ is the
population root of the FIML estimating equation. Under the stated identification
(each moment seen by some pattern) the root is unique, so the $Z$-estimator
$\hat\eta$ satisfies $\hat\eta\xrightarrow{p}\eta_0$; the argmin step of
Theorem~\ref{thm:complete} carries it to $\hat\theta$.
\end{proof}

The single load-bearing step is $\E[r_i\mid R_i=o]=0$: under MCAR, selecting a
pattern does not tilt the observed-block mean. That is exactly the step that
breaks under MAR, where selection looks at the observed values and shifts their
conditional mean; Section~\ref{sec:mar} is the accounting of that shift.

\section{MAR: the real question}\label{sec:mar}

Now let the missingness depend on the data. We use the structure of the
experiment-23 (Savalei--Bentler) design, which is also the canonical MAR setup:
a block $A$ of variables is \emph{always observed} and is the only thing the
mechanism may look at,
\[
  P(R=r\mid X) = g_r(X_A),
\]
and the remaining variables $B$ are observed in various patterns selected by
$g_r(X_A)$. Write the always-observed/conditional decomposition of the Gaussian
model.

\begin{lemma}[Variation-independent factorization]\label{lem:factor}
The Gaussian density factors as
\[
  \phi(x;\mu,\Sigma)
   = \phi(x_A;\mu_A,\Sigma_{AA})\,
     \phi\!\big(x_B;\ \alpha+\Gamma x_A,\ \Sigma_{BB\cdot A}\big),
\]
with $\Gamma=\Sigma_{BA}\Sigma_{AA}^{-1}$,
$\alpha=\mu_B-\Gamma\mu_A$, and
$\Sigma_{BB\cdot A}=\Sigma_{BB}-\Sigma_{BA}\Sigma_{AA}^{-1}\Sigma_{AB}$. The map
$(\mu,\Sigma)\leftrightarrow(\mu_A,\Sigma_{AA},\alpha,\Gamma,\Sigma_{BB\cdot A})$
is a smooth bijection with the two coordinate blocks varying independently.
Consequently the FIML log-likelihood splits exactly,
\[
  \underbrace{\sum_{i} \log\phi(x_{i,A};\mu_A,\Sigma_{AA})}_{\text{all units: }A\text{ always observed}}
  \;+\;
  \underbrace{\sum_{i:\,B\text{ obs.}} \log\phi\big(x_{i,B};\alpha+\Gamma x_{i,A},\Sigma_{BB\cdot A}\big)}_{\text{selected units: Gaussian regression of }B\text{ on }A},
\]
and the two halves are maximized separately.
\end{lemma}

The first half is a marginal model on the fully observed $A$, so by
Lemma~\ref{lem:sat} it estimates $(\mu_{A,0},\Sigma_{AA,0})$ consistently,
distribution-free. Everything hinges on the second half: a Gaussian
linear-homoscedastic regression of $B$ on $A$, fit on the \emph{selected}
complete cases.

\begin{theorem}[MAR consistency condition]\label{thm:mar}
Under MAR-on-$A$ with finite second moments and FIML identification, the
Gaussian FIML estimator is consistent for $(\mu_0,\Sigma_0)$ if and only if the
full-data conditional moments of the incomplete block satisfy
\begin{itemize}
  \item[(i)] $\E[X_B\mid X_A]=\alpha_0+\Gamma_0 X_A$ is \textbf{linear} in
        $X_A$, and
  \item[(ii)] $\Cov[X_B\mid X_A]=\Sigma_{BB\cdot A,0}$ does \textbf{not depend}
        on $X_A$ (homoscedastic).
\end{itemize}
Under (i)--(ii) consistency holds for arbitrarily non-normal marginals. If
either fails, $\hat\eta$ carries an asymptotic bias whose magnitude is set by the
selection reweighting of $X_A$.
\end{theorem}

\begin{proof}
The $A$-block is consistent by Lemma~\ref{lem:factor}. For the conditional
block, the complete cases are selected with probability $g(X_A)$ depending
\emph{only} on $X_A$. Hence the conditional law $X_B\mid X_A$ is unchanged by
selection ($f_{\mathrm{sel}}(x_B\mid x_A)=f(x_B\mid x_A)$); only the marginal of
$X_A$ is reweighted to $f_{\mathrm{sel}}(x_A)\propto g(x_A)f(x_A)$.

\emph{Mean and slope.} The Gaussian conditional MLE solves the selected normal
equations $\E_{\mathrm{sel}}\big[(1,X_A')'\,(X_B-\alpha-\Gamma X_A)'\big]=0$.
Evaluate the left side at the true full-data coefficients $(\alpha_0,\Gamma_0)$
and use the tower property over $X_B\mid X_A$ (selection-invariant):
\[
  \E_{\mathrm{sel}}\big[(1,X_A')'\,r(X_A)'\big],
  \qquad
  r(X_A) = \E[X_B\mid X_A]-\alpha_0-\Gamma_0 X_A .
\]
Under (i), $r\equiv 0$, so this vanishes \emph{for every} selection weight $g$;
$(\alpha_0,\Gamma_0)$ is therefore the population root of the selected equation
and is estimated consistently.

\emph{Conditional covariance.} The residual MLE targets
$\E_{\mathrm{sel}}\big[(X_B-\alpha_0-\Gamma_0X_A)(\cdot)'\big]
 = \E_{\mathrm{sel}}\!\big[\Cov(X_B\mid X_A)\big] + \E_{\mathrm{sel}}[\,r\,r'\,]$.
Under (i) the second term is zero, and under (ii) the first equals the constant
$\Sigma_{BB\cdot A,0}$ regardless of the $g$-weighting. So
$\Sigma_{BB\cdot A,0}$ is recovered.

Mapping back through the continuous bijection of Lemma~\ref{lem:factor}
($\mu_B=\alpha+\Gamma\mu_A$, $\Sigma_{BA}=\Gamma\Sigma_{AA}$,
$\Sigma_{BB}=\Sigma_{BB\cdot A}+\Gamma\Sigma_{AA}\Gamma'$) yields consistency for
$(\mu_0,\Sigma_0)$.

\emph{Necessity.} If (i) fails, $r\not\equiv0$, and the selected normal equations
are solved by the linear projection over $f_{\mathrm{sel}}(x_A)\propto
g(x_A)f(x_A)$, which differs from the full-population projection whenever $g$ is
non-constant: $\E_{\mathrm{sel}}[(1,X_A')'r]\ne0$ generically, so the probability
limit of $(\hat\alpha,\hat\Gamma)$ is biased. If (ii) fails,
$\E_{\mathrm{sel}}[\Cov(X_B\mid X_A)]$ is a $g$-weighted average of a
non-constant conditional covariance and differs from the unweighted
$\Sigma_{BB\cdot A,0}$. Either way $\plim\hat\eta\ne\eta_0$.
\end{proof}

\begin{corollary}[Why marginal shape is irrelevant]\label{cor:marginal}
Skewness and kurtosis of the marginal law of $X$ enter only through the
\emph{shape} of $f(x_B\mid x_A)$ and $f(x_A)$, never through conditions
(i)--(ii), which constrain only the first two conditional moments. A heavily
skewed or heavy-tailed marginal with linear, homoscedastic conditionals (the
regime that Vale--Maurelli, NORTA, and most copula constructions approximate
well) therefore gives a \textbf{consistent} FIML estimator under MAR.
\end{corollary}

\begin{remark}[The elliptical sharpening]
Elliptical non-normal laws are the cleanest cautionary example. They have
$\E[X_B\mid X_A]$ \emph{linear} (so (i) holds) but
$\Cov[X_B\mid X_A]=w(X_A)\,\Sigma_{BB\cdot A}$ \emph{heteroscedastic} (so (ii)
fails unless $w$ is constant). Under MAR on an elliptical population, FIML stays
consistent for $\mu_0$ and for the slope/regression part
$\Sigma_{BA}\Sigma_{AA}^{-1}$, but is biased for the residual covariance
$\Sigma_{BB\cdot A}$, hence for $\Sigma_{BB}$. The mean structure is robust; the
dependent-block variance is the fragile piece. ``FIML is consistent under MAR''
must be qualified at least this far.
\end{remark}

\section{Reconciliation, and what the simulations show}

\paragraph{The literature claim is correct under its hypotheses.}
Rubin's ignorability \cite{rubin1976} delivers valid likelihood inference under
MAR \emph{when the complete-data model is correctly specified}. For the Gaussian
model, ``correctly specified'' does not mean full normality; it means the first
two conditional moments are Gaussian-shaped, which is exactly (i)--(ii). That is
why the claim survives marginal non-normality. Yuan and Bentler
\cite{yuanbentler2000} and Arminger and Sobel \cite{armingersobel1990} prove
consistency and asymptotic normality of the normal-theory estimator under MAR;
their robustness is to \emph{marginal} non-normality, which enters only the
sandwich meat $J_\eta$ (the third- and fourth-moment information) and so moves
the standard errors and the test reference law, \emph{not} the location of the
estimating-equation root. Conditional misspecification, by contrast, moves the
root itself; that is the irreducible ``estimator-target'' bias the companion
note flags as something no reference-law correction (MLR, SB, FMG) can repair
\cite{savalei2010}.

So the textbook statement and the folklore are not in conflict. Read ``correct
model'' as ``correct first two conditional moments'' and both are true:
consistent under (i)--(ii), inconsistent otherwise.

\paragraph{Reading the simulations.}
Two phenomena look alike at finite $N$ and must be separated:
\begin{itemize}
  \item[(a)] \emph{Genuine asymptotic bias} (conditions (i)--(ii) fail): the
        bias converges to a nonzero constant as $N\to\infty$.
  \item[(b)] \emph{Finite-sample bias of a consistent estimator} (conditions
        hold): the bias is $O(N^{-1})$ or $O(N^{-1/2})$ and vanishes.
\end{itemize}
The experiment-23 MAR is mild on exactly the axis that matters here: selection
is a coarse \emph{sign} rule on the always-observed $x_1,x_2$, and the
Vale--Maurelli / inverse-Gaussian / power-law generators keep the
missing-given-observed conditional means close to linear and the conditional
covariances close to constant. The design therefore nearly satisfies
(i)--(ii), its true asymptotic bias is $\approx0$, and what appears in the
tables is case (b): a small-sample bias that shrinks with $N$. That is fully
consistent with the estimator being consistent.

\paragraph{How to settle it empirically.}
To decide (a) vs (b) for any given cell, regress the absolute parameter bias on
$1/N$ (or $1/\sqrt N$) across the sample-size ladder: a near-zero intercept
certifies consistency (the slope is the benign finite-sample term); a clearly
positive intercept certifies inconsistency. The test-statistic analogue is the
note's $\widehat{\mathrm{ncp}}=\texttt{mean\_base}-\texttt{mean\_trace}$, which
must shrink toward zero with $N$ under (b) and persist under (a). Predicted
outcome for experiment 23: intercepts indistinguishable from zero, i.e.\ a
finite-sample artifact, not inconsistency.

\paragraph{Verdict.}
\begin{itemize}
  \item Non-normal complete data and MCAR: consistent, unconditionally
        (Theorems~\ref{thm:complete} and \ref{thm:mcar}).
  \item MAR: consistent if and only if the incomplete-on-observed conditional
        means are linear and conditional covariances homoscedastic
        (Theorem~\ref{thm:mar}); marginal non-normality alone does not break it
        (Corollary~\ref{cor:marginal}). Your instinct is right for
        \emph{general} non-normal MAR; the paper's claim is right under the
        conditional-moment conditions, which the experiment-23 generators very
        nearly meet, so the residual MAR bias you see is finite-sample and
        should vanish with $N$.
\end{itemize}

\begin{thebibliography}{9}

\bibitem{andersonamemiya1988}
Anderson, T. W., \& Amemiya, Y. (1988).
The asymptotic normal distribution of estimators in factor analysis under
general conditions.
\emph{The Annals of Statistics}, 16(2), 759--771.
\href{https://doi.org/10.1214/aos/1176350834}{doi:10.1214/aos/1176350834}.

\bibitem{armingersobel1990}
Arminger, G., \& Sobel, M. E. (1990).
Pseudo-maximum likelihood estimation of mean and covariance structures with
missing data.
\emph{Journal of the American Statistical Association}, 85(409), 195--203.
\href{https://doi.org/10.1080/01621459.1990.10475324}{doi:10.1080/01621459.1990.10475324}.

\bibitem{browne1984}
Browne, M. W. (1984).
Asymptotically distribution-free methods for the analysis of covariance
structures.
\emph{British Journal of Mathematical and Statistical Psychology}, 37(1),
62--83.
\href{https://doi.org/10.1111/j.2044-8317.1984.tb00789.x}{doi:10.1111/j.2044-8317.1984.tb00789.x}.

\bibitem{rubin1976}
Rubin, D. B. (1976).
Inference and missing data.
\emph{Biometrika}, 63(3), 581--592.
\href{https://doi.org/10.1093/biomet/63.3.581}{doi:10.1093/biomet/63.3.581}.

\bibitem{savalei2010}
Savalei, V. (2010).
Expected versus observed information in SEM with incomplete normal and
nonnormal data.
\emph{Psychological Methods}, 15(4), 352--367.
\href{https://doi.org/10.1037/a0020143}{doi:10.1037/a0020143}.

\bibitem{white1982}
White, H. (1982).
Maximum likelihood estimation of misspecified models.
\emph{Econometrica}, 50(1), 1--25.
\href{https://doi.org/10.2307/1912526}{doi:10.2307/1912526}.

\bibitem{yuanbentler2000}
Yuan, K.-H., \& Bentler, P. M. (2000).
Three likelihood-based methods for mean and covariance structure analysis with
nonnormal missing data.
\emph{Sociological Methodology}, 30(1), 165--200.
\href{https://doi.org/10.1111/0081-1750.00078}{doi:10.1111/0081-1750.00078}.

\end{thebibliography}

\end{document}
